\documentclass[10pt]{article}
\usepackage{icmlko}

\icmlmeta{IP 파운데이션 월드모델}{243}{안쪽 눈금은 큰 것만 읽는다}{노트 241의 ``거꾸로 고른다''를 정정한다 --- 거꾸로가 아니라 그 크기 대에 눈금이 없다}

\begin{document}

\twocolumn[
\icmlheader

\begin{abstract}
노트 241은 도메인별 기울기 쌍을 안쪽에서 고르면 유보가 내려간다는 것을 보고
\textbf{``안쪽이 거꾸로 고른다''}고 적었다. 너무 셌다. 눈금을 제대로 재려고
축 열일곱 개를 \emph{하나씩} 빼서 안쪽 이득과 유보 이득을 짝지어 열일곱 점을
만들었다. \textbf{안쪽은 거꾸로가 아니다} --- 전체 spearman이 $+$0.69(F23) ·
$+$0.75(F21)이고 부호가 13/17 맞는다. 그런데 \textbf{유보 효과가 0.010 안쪽인
열두 점만 보면 spearman $+$0.25 · 부호 6/12 로 동전 던지기}다. 손 축 다섯은
유보 효과가 $-$0.013에서 $-$0.064까지 크고 다섯 다 부호가 맞는다.
\textbf{안쪽은 큰 것을 읽고 작은 것을 못 읽는다.} 노트 241의 기울기 쌍은
전부 $|\Delta| \le 0.010$ 자리였다 --- 즉 눈금이 없는 칸에서 잰 것이고,
거기서 우연히 반대로 떨어졌다. 결론은 그대로 유지된다(그 이득은 여전히 못
고른다). 바뀌는 것은 \textbf{이유}이고, 그 이유가 이 실험실의 거의 모든
후보에 걸린다 --- 우리가 다투는 변화는 대개 0.010 자리다. 곁가지로 셋을 닫았다:
얇힌 판에서 $\tau$ 칼날이 사라졌고(잡음 열 열둘에 0/13, 판 흔들림 0.0056
$\to$ 0.0025), 남은 겹말 세 쌍은 \textbf{직교화하면 더 나쁘고}($-$0.0070,
$t{=}-3.65$), \texttt{momentum} 은 유일하게 안 겹치는 시계열 축인데 빼면
$-$0.0094($t{=}-2.21$)다. 그리고 이 실험 자체가 노트 187의 검사에 한 번
걸렸다 --- \textbf{내 스크립트가 손 축을 못 빼고 있었다.}
\end{abstract}
]

\section{노트 241을 정정한다}

노트 241은 (축, 도메인) 쌍에 전용 기울기를 주는 선택을 안쪽에서 골라 봤고,
안쪽 상위 셋을 채택하면 유보가 $-$0.0070 이라는 것을 봤다. 거기서
``안쪽이 거꾸로 고른다''고 적었다. 41쌍 중 상위 셋에서 나온 한 점으로 눈금
전체를 말한 것이다.

눈금은 이렇게 잰다 --- 크기가 서로 다른 변화를 여럿 만들고, 각각에 대해
안쪽과 유보를 짝지어 본다. 제일 값싸고 균질한 변화는 \textbf{축 하나 빼기}다.
열일곱 축을 하나씩 뺐다.

\section{노트 187이 내 스크립트를 잡았다}

처음 돌린 결과에서 손 축 다섯의 유보 이득이 \textbf{정확히 $+$0.0000} 이었다.
노트 187의 제일 값싼 검사가 그대로 걸렸다 --- ``판이 소수 넷째 자리까지 안
움직이면 효과가 없는 게 아니라 모형에 안 닿은 것이다.''

닿지 않은 이유는 단순했다. 내 스크립트는 \texttt{\_trend()} ·
\texttt{\_cal()} · \texttt{\_wiki()} 가 돌려주는 dict 에서 축 이름을
걸러냈는데, \textbf{손 축 다섯은 거기서 오지 않는다} --- \texttt{*\_axes.json}
에서 \texttt{harness.load} 로 들어온다. 그래서 판이 한 글자도 안 바뀌었고
안쪽만 $-$0.04씩 내려갔다. 그 다섯 점이 상관을 통째로 음수로 끌고 갔다.

마스크로 껐다(\texttt{M[:, j] = 0}). 다시 재니 손 축은 유보에서도 크게
움직인다 --- \texttt{entry\_friction} $-$0.0637 · \texttt{target\_breadth}
$-$0.0304 · \texttt{media\_push} $-$0.0150.

\section{안쪽은 큰 것만 읽는다}

\begin{figure}[t]
\centering
\includegraphics[width=\columnwidth]{figs/scale.pdf}
\caption{축 하나씩 빼서 만든 열일곱 점. 노란 칸이 $|\Delta| \le 0.010$ 이고,
이 실험실이 다투는 변화는 대개 그 안에 있다.}
\end{figure}

\begin{center}\small
\begin{tabular}{lrrr}
\hline
 & spearman & pearson & 부호 \\
\hline
열일곱 점 전부 (F21) & $+$0.752 & $+$0.704 & 13/17 \\
열일곱 점 전부 (F23) & $+$0.691 & $+$0.717 & 11/17 \\
작은 것 열둘 (F21) & $+$0.399 & --- & 8/12 \\
\textbf{작은 것 열둘 (F23)} & $\mathbf{+0.252}$ & --- & \textbf{6/12} \\
\hline
\end{tabular}
\end{center}

손 축 다섯(유보 $-$0.013 $\sim$ $-$0.064)은 다섯 다 부호가 맞는다. 시계열 ·
달력 열둘(유보 $\pm$0.007)은 6/12 다. \textbf{안쪽은 거꾸로가 아니라, 작은
칸에서 눈금이 없다.}

노트 241의 기울기 쌍은 전부 그 작은 칸에 있었다. 결론은 안 바뀐다 --- 여전히
못 고른다. 바뀌는 것은 이유이고, 이유가 바뀌면 \emph{다음에 할 일}이 바뀐다.
``안쪽 신호가 반대다''면 안쪽을 고칠 방법을 찾아야 하지만, ``그 크기 대에
분해능이 없다''면 \textbf{안쪽 조각을 늘리거나 그 크기를 안 다투거나} 둘 중
하나다.

\section{곁가지 셋을 닫는다}

\paragraph{$\tau$ 칼날이 사라졌다.} 노트 240은 순수한 잡음 열 하나로 감쇠
$\tau$ 가 여섯 씨앗 중 셋에서 꺼지는 것을 봤다($-$0.0056). 노트 242가 축을
17로 줄인 뒤 다시 재니 \textbf{잡음 열 열둘에 0/13} --- $\tau$ 는 언제나
꺼져 있고(안정 2/5) 판 흔들림이 0.0056에서 0.0025로 줄었다. 칼날은
$\tau$ 자체가 아니라 \textbf{겹말이 실린 판}에 있었다.

\paragraph{남은 겹말은 겹친 채로 둔다.} 노트 242 뒤에 남은 세 쌍은
\texttt{level} $\sim$ \texttt{volatility} 다. 능형에게는 같은 열이니
\texttt{volatility} 를 \texttt{level} 로 설명하고 남은 것(잔차)으로 바꿔
봤다. \textbf{$-$0.0070($t{=}-3.65$)} 이다. 모형이 쓰는 것은 잔차가 아니라
\emph{같은 크기를 두 번 조금 다르게 말한 것}이다.

\paragraph{\texttt{momentum} 이 제일 값지다.} 시계열 네 축 중 유일하게 안
겹치는 축인데, 빼면 $-$0.0094($t{=}-2.21$)로 \texttt{level} 다음으로 크다.
\textbf{겹말 가드가 짚지 않은 축이 제일 값진 축이었다.}

\paragraph{남은 것}
\begin{center}\small
\begin{tabular}{ll}
\hline
갈래 & 왜 \\
\hline
안쪽 조각 & 셋을 다섯 · 여섯으로 늘리면 작은 칸이 보이나 \\
0.010 아래 & 여기를 다투지 말아야 하나. 그러면 무엇을 다투나 \\
\texttt{entry\_friction} & 빼면 $-$0.0637 --- 판에서 제일 무거운 손 축 \\
\texttt{venue\_prominence} & 안쪽 $-$0.0010 인데 유보 $-$0.0089. 왜 어긋나나 \\
\hline
\end{tabular}
\end{center}

\paragraph{규약} 안쪽 수를 근거로 쓸 때 \textbf{그 효과의 크기를 같이 적는다.}
$|\Delta| \le 0.010$ 이면 안쪽은 근거가 아니라 그냥 숫자다.

\end{document}
